\subsection{Holomorphic functional calculus in one variable}

THEOREM 5. --- Let A be a unital Banach algebra, not necessarily commutative. Let a be an element of A and z the germ of the identity function of C in a neighborhood of SpA(a). There exists a unique unital continuous morphism \(\varphi_{a}\) of \(\mathcal{O}(\mathrm{Sp}_{\mathrm{A}}(a))\) to A such that \(\varphi_{a}(z)=a\) .

The image of \(\varphi_{a}\) is contained in the full closed subalgebra of A generated by \(a\) . In particular, it is contained in the bicommutant of \(a\) .

Let us prove the existence of the morphism \(\varphi_{a}\) . Let B be the closed full subalgebra of A generated by \(a\) . It is commutative, and one has \(\mathrm{Sp}_{\mathrm{B}}(a) = \mathrm{Sp}_{\mathrm{A}}(a)\) (I, p.~5, no. 5). The mapping \(\Theta_{a}\) of holomorphic functional calculus on B is a continuous unital morphism from \(\mathcal{O}(\mathrm{Sp}_{\mathrm{B}}(a))\) to B such that \(\Theta_{a}(z) = a\) (Theorem 1 of I, p.~51). The composed morphism of \(\Theta_{a}\) and of the canonical injection of B into A is a continuous unital morphism \(\varphi_{a}\) of \(\mathcal{O}(\mathrm{Sp}_{\mathrm{A}}(a))\) into A such that the image of \(z\) is \(a\) .

Let us prove uniqueness. Let \(\varphi_{a}^{\prime}\) be a continuous unital morphism from \(\mathcal{O}(\mathrm{Sp}_{\mathrm{A}}(a))\) to A such that \(\varphi_{a}^{\prime}(z)=a\) . Then \(\varphi_{a}\) and \(\varphi_{a}^{\prime}\) coincide on the set of germs of polynomials in a neighborhood of \(\mathrm{Sp}_{\mathrm{A}}(a)\) , hence on the set of germs of holomorphic rational fractions in a neighborhood of \(\mathrm{Sp}_{\mathrm{A}}(a)\) . Now these germs are dense in \(\mathcal{O}(\mathrm{Sp}_{\mathrm{A}}(a))\) (I, p.~69, Th. 3). This implies that \(\varphi_{a}=\varphi_{a}^{\prime}\) .

The construction of \(\varphi_{a}\) shows that its image is contained in the commutative subalgebra B, which is contained in the bicommutant of a (I, p.~6).

If the radical of the algebra A is zero, the uniqueness of the morphism \(\varphi_{a}\) is valid without requiring it to be continuous (cf.~prop. 9 of I, p.~40). This is not the case in general, cf.~G. R. ALLAN, Embedding the algebra of formal power series in a Banach algebra, Proc. London Math. Soc. (3) 25 (1972), 329--340.

For every Banach algebra A, every element \(a\) of A and every germ \(f \in \mathcal{O}(\mathrm{Sp}_{\mathrm{A}}(a))\) , we denote by \(f(a)\) the element \(\varphi_{a}(f)\) of Theorem 5. If A is a commutative Banach algebra, this element \(f(a)\) coincides with the element \(f(a)\) provided by the holomorphic functional calculus on a commutative Banach algebra (Theorem 1 of I, p.~51).

Let B be the closed full subalgebra of A generated by \(a\) , so that \(\mathrm{Sp}_{\mathrm{A}}(a) = \mathrm{Sp}_{\mathrm{B}}(a)\) . The element \(f(a)\) of A belongs to B, and coincides with the element \(f(a)\) computed relative to the algebra B.

PROPOSITION 7. --- Let A and B be unital Banach algebras and \(\varphi\) a continuous unital morphism from A to B. Let \(a \in A\) . Then \(\mathrm{Sp}_{\mathrm{B}}(\varphi(a)) \subset \mathrm{Sp}_{\mathrm{A}}(a)\) and we have \(\varphi(f(a)) = f(\varphi(a))\) for all \(f \in \mathcal{O}(\mathrm{Sp}_{\mathrm{A}}(a))\) . In particular, for every \(\chi \in \mathsf{X}(\mathrm{A})\) , one has \(\chi(f(a)) = f(\chi(a))\) .

This follows from prop. 3 of I, p.~66.

PROPOSITION 8. --- Let A be a unital Banach algebra and \(a \in A\) . Let \(f \in \mathcal{O}(\mathrm{Sp}(a))\) . We have \(f(\mathrm{Sp}_{\mathrm{A}}(a)) = \mathrm{Sp}_{\mathrm{A}}(f(a))\) . Moreover, for every \(g \in \mathcal{O}(\mathrm{Sp}_{\mathrm{A}}(f(a)))\) , one has \(g \circ f \in \mathcal{O}(\mathrm{Sp}_{\mathrm{A}}(a))\) and \(g(f(a)) = (g \circ f)(a)\) .

This follows from th. 4.

PROPOSITION 9. --- Let A be a unital Banach algebra and \(a \in A\) . Let U be an open neighborhood of \(\mathrm{Sp}_{\mathrm{A}}(a)\) and \(f \in \mathcal{O}(\mathrm{U})\) . Let moreover V be a compact neighborhood of \(\mathrm{Sp}_{\mathrm{A}}(a)\) contained in U such that V is a piece of U with oriented boundary \(\partial\mathrm{V}\) .

For every integer \(n \geqslant 0\) , the map \(z \mapsto f(z)(z - a)^{-n-1}\) is continuous on \(\partial V\); the differential form \(z \mapsto f(z)(z - a)^{-n-1} dz\) is integrable on \(\partial V\), and we have

\[
f ^ {(n)} (a) = \frac {n !}{2 i \pi} \int_ {\partial \mathrm{V}} f (z) (z - a) ^ {- n - 1} d z\tag{7}
\]

where \(f^{(n)}\in\mathcal{O}(\mathrm{U})\) is the n-th derivative of \(f\) .

Proceed by induction on \(n\). When \(n = 0\) , the result follows from prop. 1 of I, p.~61. Now suppose that the assertion of the proposition is true for the integer \(n \geqslant 0\) . Let \(g \in \mathcal{O}(\mathbf{C} - \mathrm{Sp}_{\mathrm{A}}(a); \mathrm{A})\) be the holomorphic function defined by \(g(z) = (z - a)^{-n-1} f(z)\). The differential form \(g'(z) dz = dg\) is of class \(\mathrm{C}^1\) ; since the piece V is compact, Stokes' formula (VAR, R2, p.~47, 11.2.3) implies

\[
\int_ {\partial \mathrm{V}} g ^ {\prime} (z) d z = \int_ {\partial \mathrm{V}} d g = 0.
\]

Since \(g'(z) = (z - a)^{-n-1} f'(z) - (n+1)(z - a)^{-n-2} f(z)\) , from which one deduces

\[
\int_ {\partial \mathrm{V}} f ^ {\prime} (z) (z - a) ^ {- n - 1} d z = (n + 1) \int_ {\partial \mathrm{V}} f (z) (z - a) ^ {- n - 2} d z.
\]

Applying the induction hypothesis to \(f'\) , one therefore obtains

\[
\frac {2 i \pi}{n !} f ^ {(n + 1)} (a) = (n + 1) \int_ {\partial \mathrm{V}} f (z) (z - a) ^ {- n - 2} d z,
\]

which is the assertion of the proposition for the integer \(n+1\) . This concludes the proof.

PROPOSITION 10. --- Let A be a unital Banach algebra and U an open subset of C.

\begin{enumerate}
\def\labelenumi{\alph{enumi})}
\tightlist
\item
  The set \(\Omega\) of \(a \in \mathrm{A}\) such that \(\operatorname{Sp}_{\mathrm{A}}(a) \subset \mathrm{U}\) is open in A;
\item
  Let \(f \in \mathcal{O}(\mathrm{U})\) . The map \(a \mapsto f(a)\) of \(\Omega\) in A is holomorphic, and in particular continuous.
\end{enumerate}

Let \(a \in \Omega\) . There exists a compact neighborhood V of \(\mathrm{Sp}_{\mathrm{A}}(a)\) contained in U which is a piece of U (VAR, R2, p.~46 and p.~47, 11.1.3, d)).

Since the resolvent of \(a\) tends to 0 at infinity (th. 1, \(c\) ) of I, p.~24), the map \(z \mapsto \| (z - a)^{-1}\|\) is bounded on \(\mathbf{C} - \mathring{\mathrm{V}}\) . Let M denote its upper bound. If \(h \in \mathrm{A}\) is such that \(\| h\| \leqslant (2\mathrm{M})^{-1}\) and if \(z \in \mathbf{C} - \mathring{\mathrm{V}}\) , one has

\[
z - (a + h) = (1 - h (z - a) ^ {- 1}) (z - a)
\]

and \(\|h(z-a)^{-1}\|\leqslant\frac{1}{2}\) , hence \(z-(a+h)\) is invertible and its inverse satisfies

\[
(z - (a + h)) ^ {- 1} = (z - a) ^ {- 1} \sum_ {n = 0} ^ {\infty} (h (z - a) ^ {- 1}) ^ {n}\tag{8}
\]

with \(\| (h(z - a)^{-1})^n\| \leqslant 2^{-n}\) (prop. 2 of I, p.~22). Thus, \(\mathrm{Sp}_{\mathrm{A}}(a + h)\) is contained in V, hence in U, which proves that \(\Omega\) is open in A.

Let \(f \in \mathcal{O}(\mathrm{U})\) . Denote by \(m\) the supremum of \(|f(z)|\) for \(z \in \partial \mathrm{V}\) . Let \(a \in \mathrm{A}\) . For every \(h \in \mathrm{A}\) such that \(\| h \| \leqslant (2\mathrm{M})^{-1}\) , one has

\[
f (a + h) = \frac {1}{2 i \pi} \int_ {\partial \mathrm{V}} f (z) (z - (a + h)) ^ {- 1} d z
\]

(prop. 9). The series (8) converges uniformly on the boundary of V, hence

\[
f (a + h) = \sum_ {n = 0} ^ {+ \infty} f _ {a, n} (h)
\]

where the map \(f_{a,n}\) from A to A is defined by

\[
f _ {a, n} (h) = \frac {1}{2 i \pi} \int_ {\partial \mathrm{V}} f (z) (z - a) ^ {- 1} (h (z - a) ^ {- 1}) ^ {n} d z.
\]

For every \(n \in N\) , the function \(f_{a,n}\) is a continuous homogeneous polynomial function of degree n. Moreover, it follows that

\[
\left\| f _ {a, n} (h) \right\| \leqslant \frac {m M}{\pi} \left(\int_ {\partial V} \| d z \|\right) 2 ^ {- (n + 1)}
\]

(INT, VI, §2, no. 3, prop. 5). The series \(\sum_{n} f_{a,n}(h)\) is therefore absolutely convergent for \(\|h\| \leqslant (2M)^{-1}\). This proves that the map which to \(a\) , associates \(f(a)\) is holomorphic on \(\Omega\) (VAR, R1, p.~26, 3.2.1).

PROPOSITION 11. --- Let A be a unital Banach algebra, \(a \in A\) and U an open neighborhood of \(\mathrm{Sp}_{\mathrm{A}}(a)\) . Denote by \(\delta\) the distance from \(\mathrm{Sp}_{\mathrm{A}}(a)\) to \(\mathbf{C} - \mathbf{U}\) . Let \(f \in \mathcal{O}(\mathbf{U})\) .

\begin{enumerate}
\def\labelenumi{\alph{enumi})}
\item
  For every real number η such that 0 \textless{} η \textless{} δ, there exists a real number C ≥ 0 such that \(\|f^{(n)}(a)\| \leqslant \mathrm{Cn}!\eta^{-n}\) for every integer \(n \in N\) ;
\item
  If \(b \in A\) commutes with \(a\) and if \(\varrho(b) < \delta\) , one has \(\mathrm{Sp}_{\mathrm{A}}(a + b) \subset \mathrm{U}\) , and
\end{enumerate}

\[
f (a + b) = \sum_ {n = 0} ^ {\infty} \frac {f ^ {(n)} (a)}{n !} b ^ {n},
\]

where the series converges absolutely.

Let \(\eta\) be a real number such that \(0 < \eta < \delta\) . Denote by \(\varepsilon = \delta - \eta > 0\) . Let K be the compact neighborhood of \(\mathrm{Sp}_{\mathrm{A}}(a)\) consisting of the points of C whose distance to \(\mathrm{Sp}_{\mathrm{A}}(a)\) is \(\leqslant \varepsilon / 2\) . Since \(f\) is holomorphic in every open disk of radius \(\eta + \varepsilon / 2\) whose center belongs to K, there exists, by the Cauchy inequalities (VAR, R1, p.~29, 3.3.4), a real number \(C \geqslant 0\) such that

\[
\sup _ {z \in K} \frac {| f ^ {(n)} (z) |}{n !} \leqslant \frac {C}{\eta^ {n}}
\]

for every integer \(n \geqslant 0\) . Then assertion a) follows from prop. 1 of I, p.~61 applied to \(f^{(n)}\) and to a piece V contained in K.

Let \(b\) be an element of A commuting with \(a\) such that \(\varrho(b) < \delta\) . By replacing A by the closed full subalgebra B generated by \(a\) and \(b\) , which satisfies \(\mathrm{Sp}_{\mathrm{A}}(a) = \mathrm{Sp}_{\mathrm{B}}(a)\) and \(\mathrm{Sp}_{\mathrm{A}}(a + b) = \mathrm{Sp}_{\mathrm{B}}(a + b)\) , one reduces, in order to prove \(b\) ), to the case where A is commutative.

Since \(\varrho(b) < \delta\) , one can choose \(\eta\) such that \(\varrho(b) < \eta < \delta\) . Let \(V_1\) be the set of points of \(\mathbf{C}\) whose distance to \(\mathrm{Sp}_{\mathrm{A}}(a)\) is \(< \delta - \eta\) , and \(V_2\) be the open disk with center 0 and radius \(\eta\) in \(\mathbf{C}\) . Let \(g\) be the map \((z_1, z_2) \mapsto z_1 + z_2\) from \(V_1 \times V_2\) into U. Then \(h = f \circ g\) is the map \((z_1, z_2) \mapsto f(z_1 + z_2)\) from \(V_1 \times V_2\) into \(\mathbf{C}\) . We have \(\mathrm{Sp}_{\mathrm{A}}^{2}(a, b) \subset V_1 \times V_2\) , hence \(\mathrm{Sp}_{\mathrm{A}}(a + b) \subset U\) (cf. Cor. 2 of I, p. 67), and moreover \(f(a + b) = h(a, b)\) . By th. 4 of I, p.~72. Now, in the space \(\mathcal{O}(V_1 \times V_2)\) , one has

\[
h (z _ {1}, z _ {2}) = \sum_ {n \geqslant 0} \frac {f ^ {(n)} (z _ {1})}{n !} z _ {2} ^ {n},
\]

therefore the series

\[
\sum_ {n \geqslant 0} \frac {f ^ {(n)} (a)}{n !} b ^ {n}
\]

converges in A and its sum is \(h(a,b)=f(a+b)\). Moreover, this series is absolutely convergent by assertion (a).

